% !TEX program = xelatex
% Драфт для ПЗ-3. Класс оформления сохранён из шаблона преподавателя.
\documentclass{misisarticle}
\usepackage{tabularx}
\usepackage{needspace}
\grnti{28.23.15, 28.23.37}
\asjc{1707, 1702}
\oecd{1.02}
\title{Исследование методов доменной адаптации нейронных сетей для семантической сегментации при переносе с синтетических изображений на реальные}
\misisauthor{Н.~М.~Коршунов}{кафедра инженерной кибернетики НИТУ «МИСиС»}{Москва, Россия}{m2303488@edu.misis.ru}
\abstract{Рассматривается перенос моделей семантической сегментации с размеченных синтетических изображений на реальные изображения без доступной обучающей разметки целевого домена. Цель работы --- систематизировать математические механизмы доменной адаптации и определить условия их корректного сравнения. Сопоставляются статистическое и транспортное согласование, состязательное обучение, минимизация энтропии, самообучение, согласованность предсказаний и контрастивные критерии. На примере опубликованных исследований показано, почему улучшение всей адаптационной схемы следует отличать от эффекта отдельного компонента. Сформулирована постановка переноса GTA5~$\rightarrow$~Cityscapes и предложен протокол будущего сравнения по mIoU, качеству отдельных классов и вычислительным затратам. Работа представляет обзор и проект исследования; собственные экспериментальные результаты на данном этапе не приводятся.}
\keywords{Доменная адаптация; перенос с синтетических данных на реальные (Synthetic-to-Real); семантическая сегментация; сдвиг домена (domain shift); обучение без разметки целевого домена (UDA); псевдоразметка (pseudo-labeling)}
\begin{document}
\maketitle
\section{Введение}
Семантическая сегментация сопоставляет каждому пикселю изображения класс объекта: дорога, автомобиль, пешеход и другие категории. Пиксельная разметка реальных сцен требует значительных трудозатрат. Синтетические сцены позволяют получать изображения и соответствующие метки из виртуальной среды, однако различия текстур, освещения, геометрии и частот классов затрудняют перенос модели на фотографии.

Доменная адаптация без целевой разметки (unsupervised domain adaptation, UDA) использует размеченные изображения исходного домена и неразмеченные изображения целевого домена. Важен не только выбор сегментационной сети, но и критерий, позволяющий извлечь обучающий сигнал из целевых данных. В AdaptSegNet используется состязательное согласование карт предсказаний~\cite{tsai}; DACS сочетает псевдоразметку с междоменным смешиванием изображений~\cite{dacs}; DAFormer показывает существенное влияние архитектуры и стратегии обучения~\cite{daformer}.

Итоговые метрики этих работ не позволяют ранжировать отдельные математические принципы: одновременно меняются модель, начальные веса, аугментации и правила обучения. Цель работы --- сопоставить механизмы адаптации и сформулировать условия их проверки. Аналитический результат включает систематизацию критериев, их ограничений и правил сравнения опубликованных оценок.

\section{Анализ работ}
\subsection{Смежные исследования и исходная проблема}
В студенческой работе С.~Аскари Хеммат рассматривается подготовка синтетических сцен с мебелью в Unreal Engine и обучение моделей на этих данных~\cite{askari}. Она связана с настоящей темой способом получения изображений, но приведённые в ней показатели детекции не являются оценкой UDA на дорожных сценах. Л.~Б.~Комендала рассматривает сегментацию дорожных изображений CamVid~\cite{komendala}; эта работа близка по прикладной задаче, однако не задаёт перенос GTA5~$\rightarrow$~Cityscapes. В исследовании И.~И.~Курочкина и соавторов нейросети применяются для классификации микроскопических изображений криогелей~\cite{kurochkin}. Это пример анализа изображений в другой предметной области; его результаты не переносятся на метрики семантической сегментации.

\subsection{Основные направления адаптации}
Статистические критерии сравнивают распределения признаков через ядерные средние или ковариации~\cite{mmd,coral}. Оптимальный транспорт строит соответствия между наблюдениями, минимизируя стоимость переноса массы~\cite{sinkhorn}. Состязательные методы обучают модель скрывать различие доменов от дополнительного классификатора. В AdaptSegNet согласование проводится на уровне пространственно организованных предсказаний сегментатора~\cite{tsai}.

Другая группа подходов использует структуру ответов на целевых изображениях. Энтропийный критерий поощряет определённые предсказания~\cite{advent}. Самообучение превращает часть предсказаний в обучающие метки. Согласованность требует устойчивости ответов к допустимым преобразованиям; в MIC ученик предсказывает классы по частично закрытому изображению, используя псевдометки учителя~\cite{mic}. Контрастивные методы, например SePiCo, добавляют семантические ограничения на расположение признаков~\cite{sepico}.

Направления пересекаются. В DACS псевдометки используются вместе со смешиванием исходных и целевых изображений и схемой «учитель --- ученик». DAFormer дополняет самообучение выборкой редких классов, регуляризацией признаков и разогревом скорости обучения. Поэтому название алгоритма обозначает целую реализацию, а не единственный математический критерий. В работе рассматриваются шесть перечисленных направлений; преобразование внешнего вида изображений не выделяется в самостоятельный предмет сравнения.

\subsection{Что показывают опубликованные оценки}
Для оценки всей схемы нужна модель, обученная только на исходном домене (source only). В AdaptSegNet многоуровневое согласование повышает mIoU с 36,6 до 42,4, то есть на 5,8 процентного пункта (п.~п.)~\cite[табл.~1]{tsai}. В DACS соответствующая пара равна 32,85 и 52,14, прирост составляет 19,29~п.~п.~\cite[табл.~1]{dacs}. Обе оценки относятся к GTA5~$\rightarrow$~Cityscapes и 19 классам, но условия обучения различаются. Разница приростов не доказывает преимущество самообучения над состязательным критерием в одинаковой модели.

Внутреннее сравнение DACS (табл.~\ref{tab:dacs}) показывает риск ухудшения при псевдоразметке. Naive Mixing смешивает целевые изображения между собой, DACS --- исходные и целевые. Различие результатов характеризует эти реализации, а не направление самообучения в целом.

\begin{table}[htbp]
\caption{Опубликованное сравнение вариантов DACS на GTA5~$\rightarrow$~Cityscapes~\cite[табл.~3]{dacs}}
\label{tab:dacs}
\centering
\setlength{\tabcolsep}{3pt}
\begin{tabularx}{\columnwidth}{@{}Xrr@{}}
\toprule
Вариант & mIoU, \% & $\Delta$, п.~п.\\
\midrule
Только исходный домен & 32,85 & ---\\
Псевдоразметка без смешивания & 22,97 & $-9{,}88$\\
Naive Mixing & 35,08 & $+2{,}23$\\
DACS & 52,14 & $+19{,}29$\\
\bottomrule
\end{tabularx}
\par\vspace{3pt}
{\footnotesize $\Delta$ рассчитана относительно 32,85. Это результаты авторов DACS, не собственные эксперименты.}
\end{table}

Эффект отдельного компонента оценивают относительно уже адаптирующейся системы. Добавление MIC к DAFormer изменяет mIoU с 68,3 до 70,6, то есть на 2,3~п.~п.~\cite[табл.~7]{mic}. Это дополнительный эффект MIC; исходные 68,3~mIoU~\cite{daformer} характеризуют сочетание архитектуры и правил обучения. Сопоставимого численного примера для CORAL в выбранном сценарии не приводится.

\section{Постановка задачи}
Пусть $\mathcal D_s=\{(x_i^s,y_i^s)\}_{i=1}^{N_s}$ --- размеченная исходная выборка, а $\mathcal D_t=\{x_j^t\}_{j=1}^{N_t}$ --- целевая выборка без доступных обучающих меток. Пространство из $C$ классов считается общим. Сеть с параметрами $\theta$ выдаёт в пикселе $u$ вектор вероятностей $p_\theta(x)_u$; итоговая метка выбирается по максимальной вероятности. Распределения изображений двух доменов различаются.

Схема разделения обучающих и проверочных данных приведена на рис.~\ref{fig:uda-protocol}.

\begin{figure}[htbp]
\centering
% Авторская схема. Все подписи набираются шрифтом статьи.
\begingroup
\setlength{\unitlength}{1mm}
\fontsize{8.5}{10}\selectfont
\linethickness{0.4pt}
\begin{picture}(86,61)
  \put(0,49){\framebox(40,12){\shortstack{GTA5: $(x^s,y^s)$\\изображения и маски}}}
  \put(46,49){\framebox(40,12){\shortstack{Cityscapes train: $x^t$\\без обучающих масок}}}
  \put(20,49){\line(0,-1){6}}
  \put(66,49){\line(0,-1){6}}
  \put(20,43){\line(1,0){46}}
  \put(43,43){\vector(0,-1){6}}
  \put(7,25){\framebox(72,12){\shortstack{Обучение общей модели\\$\min_\theta\;L_s+\sum_k\lambda_k L_k^{\mathrm{adapt}}$}}}
  \put(43,25){\vector(0,-1){11}}
  \put(45,19){\makebox(0,0)[l]{параметры $\theta$}}
  \put(0,0){\framebox(25,14){\shortstack{Cityscapes val\\изображения}}}
  \put(25,7){\vector(1,0){6}}
  \put(31,0){\framebox(24,14){\shortstack{Модель\\предсказания}}}
  \put(55,7){\vector(1,0){6}}
  \put(61,0){\framebox(25,14){\shortstack{Сравнение\\с масками val\\mIoU, IoU}}}
\end{picture}
\endgroup

\caption{Постановка GTA5~$\rightarrow$~Cityscapes: обучение без целевых меток и оценка на отдельной проверочной выборке.}
\label{fig:uda-protocol}
\end{figure}

На исходных данных минимизируется перекрёстная энтропия:
\begin{equation}
 L_s(\theta)=-\frac1{N_s}\sum_{i=1}^{N_s}\frac1{|\Omega_i|}
 \sum_{u\in\Omega_i}\log p_{\theta,y^s_{iu}}(x_i^s)_u,
 \label{eq:source}
\end{equation}
где $\Omega_i$ --- множество учитываемых пикселей. Для целевых обучающих изображений такая потеря недоступна. Общая схема вводит дополнительный критерий:
\begin{equation}
 \min_\theta\; L_s(\theta)+\sum_k\lambda_k L_k^{\mathrm{adapt}}(\theta),
 \qquad\lambda_k\geq0.
 \label{eq:total}
\end{equation}
Формула задаёт способ объединения критериев, но не требует одновременного использования всех направлений. Для состязательной схемы добавляется максимизация по параметрам дискриминатора, для транспорта --- внутренняя задача определения соответствий.

Основной сценарий --- GTA5~\cite{gta} как источник синтетических изображений и Cityscapes~\cite{cityscapes} как целевой домен. Целевые обучающие маски не используются при оптимизации. Маски отдельной проверочной выборки нужны для оценки. Если они применяются для выбора гиперпараметров или момента остановки, это должно явно отражаться в протоколе: такой режим отличается от выбора модели без целевой разметки. SYNTHIA может рассматриваться позднее как дополнительный источник; результаты для 13, 16 и 19 классов нельзя смешивать в одной оценке.

Качество оценивается средней по классам мерой пересечения с объединением:
\begin{equation}
 \mathrm{mIoU}=\frac{100}{C}\sum_{c=1}^{C}
 \frac{\mathrm{TP}_c}{\mathrm{TP}_c+\mathrm{FP}_c+\mathrm{FN}_c}.
 \label{eq:miou}
\end{equation}
Числа истинно положительных, ложноположительных и ложноотрицательных пикселей накапливаются по всей проверочной выборке. Дополнительно рассматриваются IoU каждого класса, время обучения и пиковая память ускорителя. Средняя метрика может скрывать ухудшение редких категорий, поэтому одного mIoU недостаточно.

\section{Описание методов}
\subsection{Статистическое и транспортное согласование}
Пусть $z$ --- вектор признаков, извлечённый из пространственной позиции карты сети, а $P_s$ и $P_t$ --- распределения исходных и целевых признаков. Ядерный критерий MMD сравнивает их средние вложения в пространство $\mathcal H$~\cite{mmd}:
\begin{equation}
 L_{\mathrm{MMD}}=\left\|\mathbb E_{P_s}\phi(z)-\mathbb E_{P_t}\phi(z)\right\|_{\mathcal H}^2.
\end{equation}
На практике используются значения ядра $k(z,z')=\langle\phi(z),\phi(z')\rangle$. Сеть обновляется по исходной ошибке и оценке MMD, вычисленной на выбранных признаках двух доменов. Выбор ядра, его масштаба и способа выборки пикселей влияет на результат. Полная попарная оценка требует квадратичного числа операций по числу выбранных признаков.

Deep CORAL согласует ковариации $C_s$ и $C_t$ признаков размерности $d$~\cite{coral}:
\begin{equation}
 L_{\mathrm{CORAL}}=\frac{1}{4d^2}\|C_s-C_t\|_F^2.
\end{equation}
Ковариации вычисляются после центрирования признаков. Их совпадение не гарантирует совпадения средних или полных распределений. В сегментационной модели оба критерия можно применять к выборке пространственных признаков; согласование без учёта классов способно сблизить семантически разные объекты.

Оптимальный транспорт задаёт матрицу неотрицательных переносимых масс $\gamma$. Для исходных весов $a$ и целевых весов $b$, каждый из которых суммируется в единицу, рассматривается задача
\begin{equation}
 \min_{\gamma\in\Pi(a,b)}\;
 \langle\gamma,M\rangle+\varepsilon\sum_{i,j}\gamma_{ij}(\log\gamma_{ij}-1),
\end{equation}
где $\Pi(a,b)=\{\gamma\geq0:\gamma\mathbf1=a,\gamma^T\mathbf1=b\}$, $M$ --- стоимость сопоставления, $\varepsilon>0$. Принимается $0\log0=0$. Алгоритм Синкхорна попеременно нормирует строки и столбцы ядра $K_{ij}=\exp(-M_{ij}/\varepsilon)$~\cite{sinkhorn}. После расчёта плана обновляются параметры сети. Семантическая информация может включаться в стоимость; память для матрицы попарных стоимостей ограничивает число используемых признаков.

\subsection{Состязательное согласование}
Дискриминатор $D_\varphi$ оценивает вероятность принадлежности исходному домену по признакам либо предсказаниям $r_\theta(x)$. Базовая постановка имеет вид
\begin{equation}
 \begin{split}
 V(\theta,\varphi)={}&\mathbb E_s\log D_\varphi(r_\theta(x))\\
 &+\mathbb E_t\log[1-D_\varphi(r_\theta(x))],\\
 &\min_\theta\max_\varphi\;L_s(\theta)+\lambda V(\theta,\varphi).
 \end{split}
\end{equation}
Максимизация обучает различать домены, минимизация по $\theta$ затрудняет это различение. Используются попеременные градиентные шаги; конкретные алгоритмы могут менять генераторную потерю для устойчивости. AdaptSegNet применяет этот принцип к картам классов~\cite{tsai}. Основные ограничения --- неустойчивость минимаксного обучения и возможность подавления полезных различий. Низкая точность дискриминатора сама по себе не подтверждает правильность сегментации.

\subsection{Минимизация энтропии}
Для вероятностей классов $p_c$ нормированная энтропия равна
\begin{equation}
 H(p)=-\frac1{\log C}\sum_{c=1}^{C}p_c\log p_c.
\end{equation}
Её среднее по целевым пикселям добавляется к $L_s$. Критерий стимулирует уверенные ответы, но уверенная ошибка также имеет малую энтропию. Поэтому важно контролировать вес регуляризации и распределение предсказываемых классов. В ADVENT различаются прямая минимизация энтропии MinEnt и состязательная обработка карт самоинформации AdvEnt~\cite{advent}; их нельзя считать одной и той же процедурой оптимизации.

\subsection{Самообучение с псевдоразметкой}
Модель с фиксированными на текущем шаге параметрами $\bar\theta$ формирует метки $\hat y_u=\arg\max_c p_{\bar\theta,c}(x^t)_u$. Для отбора можно использовать маску $w_u=\mathbf1[\max_c p_{\bar\theta,c}(x^t)_u\geq\tau_{\hat y_u}]$ и потерю
\begin{equation}
 L_{\mathrm{pl}}=-\frac{\sum_u w_u\log p_{\theta,\hat y_u}(x^t)_u}
 {\max(1,\sum_u w_u)}.
\end{equation}
Псевдометки и маска фиксируются при обновлении ученика. Затем предсказания уточняются, и обучение повторяется. Пороги могут зависеть от класса, поскольку единый высокий порог часто отбирает преимущественно простые и частые категории. Это обобщённая схема отбора, а не точное воспроизведение DACS или DAFormer.

В DACS части исходных изображений вместе с истинными метками смешиваются с целевыми областями и псевдометками~\cite{dacs}. DAFormer использует выборку исходных изображений с редкими классами~\cite{daformer}. Эти механизмы помогают сохранять классовое разнообразие, но не устраняют риск закрепления ошибок автоматически. Сравнение в табл.~\ref{tab:dacs} показывает, почему наличие псевдометок само по себе не гарантирует улучшения.

\subsection{Согласованность предсказаний}
Для двух допустимых преобразований одного изображения минимизируется расхождение ответов, например дивергенция Кульбака--Лейблера между ответом учителя $q_u$ и ученика $p_u$. Геометрические преобразования требуют предварительного совмещения координат. Параметры учителя могут обновляться экспоненциальным усреднением:
\begin{equation}
 \bar\theta_k=\alpha\bar\theta_{k-1}+(1-\alpha)\theta_k,\qquad 0\leq\alpha<1.
\end{equation}
Градиент через целевой ответ учителя не проводится. В MIC применяется перекрёстная энтропия между предсказаниями ученика по маскированному изображению и псевдометками учителя по полному изображению~\cite{mic}. Учитель сглаживает изменения параметров, но не получает независимой правильной разметки. Два одинаковых ответа могут быть одновременно ошибочными.

\subsection{Контрастивные критерии}
Контрастивное обучение сближает признаки одного класса и разделяет признаки разных классов. Положительные соответствия задаются исходными метками или целевыми псевдометками; отрицательные --- другими классами. Вариант с прототипами сравнивает признак с представительными векторами классов через softmax от сходства, делённого на температуру. Критерии типа InfoNCE задают относительное предпочтение положительных соответствий отрицательным.

В SePiCo классы представляются прототипами, банками признаков или распределениями~\cite{sepico}. Для исследования важны надёжность псевдометок, обновление прототипов, температура и состав отрицательных примеров. Ошибочная метка создаёт неверное семантическое соответствие. Поэтому контрастивный член следует оценивать как дополнение к фиксированной базовой схеме, отдельно проверяя его влияние на редкие классы.

\subsection{Проект сравнительного исследования}
Будущее сравнение предполагает единые разбиения данных, набор из 19 классов, начальные веса, разрешение входа и бюджет обучения. Сначала оценивается source only. Затем в той же модели по отдельности включаются выбранные критерии. Комбинации исследуются через исключение или добавление одного компонента при сохранении остальных условий. Полные реализации AdaptSegNet, DACS и DAFormer рассматриваются отдельно как ориентиры из литературы.

Для каждого запуска следует сохранять конфигурацию, случайное зерно, версию реализации, правило выбора контрольной точки, mIoU и IoU классов, время обучения и пиковую память. Если ресурсы позволяют, выполняются повторные запуски и сообщаются среднее и стандартное отклонение. Разности малой величины без оценки вариативности не считаются достаточным основанием для вывода о превосходстве метода. Снижение метрики относительно той же исходной модели фиксируется как отрицательный перенос.

Предстоит проверить, какие критерии дают прирост при одинаковых условиях и сохраняется ли он для редких классов. Набор вариантов и бюджет запусков зависят от вычислительных ресурсов. Обучение по этому протоколу ещё не проводилось.


\Needspace{12\baselineskip}
\begin{thebibliography}{99}
\bibitem{tsai} Tsai Y.-H., Hung W.-C., Schulter S. et al. Learning to Adapt Structured Output Space for Semantic Segmentation // CVPR. 2018. P.~7472--7481. \href{https://openaccess.thecvf.com/content_cvpr_2018/html/Tsai_Learning_to_Adapt_CVPR_2018_paper.html}{Текст публикации}.
\bibitem{dacs} Tranheden W., Olsson V., Pinto J., Svensson L. DACS: Domain Adaptation via Cross-Domain Mixed Sampling // WACV. 2021. P.~1379--1389. \href{https://openaccess.thecvf.com/content/WACV2021/html/Tranheden_DACS_Domain_Adaptation_via_Cross-Domain_Mixed_Sampling_WACV_2021_paper.html}{Текст публикации}.
\bibitem{daformer} Hoyer L., Dai D., Van Gool L. DAFormer: Improving Network Architectures and Training Strategies for Domain-Adaptive Semantic Segmentation // CVPR. 2022. P.~9924--9935. \href{https://openaccess.thecvf.com/content/CVPR2022/html/Hoyer_DAFormer_Improving_Network_Architectures_and_Training_Strategies_for_Domain-Adaptive_Semantic_CVPR_2022_paper.html}{Текст публикации}.
\bibitem{askari} Аскари Хеммат С. Применение синтетических данных из UnrealEngine для обучения модели сегментации мебели // Искусственный интеллект в промышленных, коммерческих, медицинских и финансовых приложениях. Вып.~4. М.: НИТУ «МИСИС», 2025. С.~13--16. \href{https://sadekov.su/Articles/kik_2025_1.pdf#page=14}{Сборник, с.~13}.
\bibitem{komendala} Комендала Л.~Б. Семантическая сегментация дорожных сцен с использованием глубокого обучения // Искусственный интеллект в промышленных, коммерческих, медицинских и финансовых приложениях. Вып.~5. М.: НИТУ «МИСИС», 2026. С.~82--86. \href{https://sadekov.su/Articles/kik_2025_2.pdf#page=83}{Сборник, с.~82}.
\bibitem{kurochkin} Kurochkin I.~I., Kurochkin I.~N., Kolosova O.~Yu., Lozinsky V.~I. Cryostructuring of Polymeric Systems: Application of Deep Neural Networks for the Classification of Structural Features Peculiar to Macroporous Poly(vinyl alcohol) Cryogels Prepared without and with the Additives of Chaotropes or Kosmotropes // Molecules. 2020. Vol.~25, no.~19. Art.~4480. \href{https://doi.org/10.3390/molecules25194480}{DOI: 10.3390/molecules25194480}.
\bibitem{mmd} Gretton A., Borgwardt K.~M., Rasch M.~J. et al. A Kernel Two-Sample Test // JMLR. 2012. Vol.~13. P.~723--773. \href{https://www.jmlr.org/papers/v13/gretton12a.html}{Текст публикации}.
\bibitem{coral} Sun B., Saenko K. Deep CORAL: Correlation Alignment for Deep Domain Adaptation. 2016. \href{https://arxiv.org/abs/1607.01719}{arXiv:1607.01719}.
\bibitem{sinkhorn} Cuturi M. Sinkhorn Distances: Lightspeed Computation of Optimal Transport // NeurIPS. 2013. Vol.~26. \href{https://arxiv.org/abs/1306.0895}{arXiv:1306.0895}.
\bibitem{advent} Vu T.-H., Jain H., Bucher M. et al. ADVENT: Adversarial Entropy Minimization for Domain Adaptation in Semantic Segmentation // CVPR. 2019. \href{https://arxiv.org/abs/1811.12833}{arXiv:1811.12833}.
\bibitem{mic} Hoyer L., Dai D., Wang H., Van Gool L. MIC: Masked Image Consistency for Context-Enhanced Domain Adaptation // CVPR. 2023. P.~11721--11732. \href{https://arxiv.org/abs/2212.01322}{arXiv:2212.01322}.
\bibitem{sepico} Xie B., Li S., Li M. et al. SePiCo: Semantic-Guided Pixel Contrast for Domain Adaptive Semantic Segmentation // IEEE TPAMI. 2023. \href{https://arxiv.org/abs/2204.08808}{arXiv:2204.08808}.
\bibitem{gta} Richter S.~R., Vineet V., Roth S., Koltun V. Playing for Data: Ground Truth from Computer Games // ECCV. 2016. \href{https://arxiv.org/abs/1608.02192}{arXiv:1608.02192}.
\bibitem{cityscapes} Cordts M., Omran M., Ramos S. et al. The Cityscapes Dataset for Semantic Urban Scene Understanding // CVPR. 2016. P.~3213--3223. \href{https://www.cityscapes-dataset.com/}{Сайт набора данных}.
\end{thebibliography}

\end{document}
